\section{Related Work}

Since the initial observation of grokking, several studies have characterized its underlying mechanisms from complementary angles.

\paragraph{Mechanistic Interpretability \cite{nanda2023progressmeasuresgrokkingmechanistic}.}
Nanda et al.~\cite{nanda2023progressmeasuresgrokkingmechanistic} showed that grokking is driven by a smooth process of internal circuit formation. Networks gradually develop generalizing circuits (e.g., discrete Fourier transforms) while simultaneously relying on memorization circuits; as training progresses, the generalizing circuits eventually dominate.\looseness=-1

\paragraph{Intrinsic Task Symmetry \cite{hwang2026intrinsic}.}
Hwang and Park~\cite{hwang2026intrinsic} propose that intrinsic task symmetries drive a three-stage training dynamic --- memorization, symmetry acquisition, and geometric organization --- where representations snap into a task-aligned geometry during generalization.\looseness=-1

Our work introduces TDA as a complementary lens: rather than isolating circuits or algebraic symmetries, we track global topological invariants to characterize how the full representation manifold evolves across the generalization transition. To our knowledge, ours is the first work to combine persistent homology with intrinsic dimension estimation in a dynamic, training-time analysis of grokking. Table~\ref{tab:related} summarizes our positioning.\looseness=-1

\begin{table}[h]
\vspace{-6pt}
\centering
\caption{Comparison of related approaches. ``Dynamic'' = analysis over
training; ``Static'' = analysis of a fully trained model.
PH = Persistent Homology.}
\label{tab:related}
\small
\begin{tabular*}{\textwidth}{@{\extracolsep{\fill}}lcccc}
\toprule
\textbf{Work} & \textbf{Analysis type} & \textbf{Temporal}
              & \textbf{Grokking} & \textbf{PH} \\
\midrule
Nanda et al.~\cite{nanda2023progressmeasuresgrokkingmechanistic}
    & Mechanistic   & Dynamic & \checkmark & \\
Hwang \& Park~\cite{hwang2026intrinsic}
    & Geometric     & Dynamic & \checkmark & \\
Naitzat et al.~\cite{naitzat2020topology}
    & Topological   & Static  &            & \checkmark \\
Ansuini et al.~\cite{ansuini2019intrinsic}
    & Geometric     & Static  &            & \\
\midrule
\textbf{This work}
    & Topological   & Dynamic & \checkmark & \checkmark \\
\bottomrule
\end{tabular*}
\vspace{-6pt}
\end{table}